\documentclass[../../../../main.tex]{subfiles}
\begin{document}

\begin{flushright}
\footnotesize\textit{【自動文字起こし・要確認】}
\end{flushright}

{\bf [解]}
\begin{align*}
a_1&=1，\quad S_n\xrightarrow{n\to\infty}1\quad\cdots①\\
n(n-2)a_{n+1}&=S_n\quad\cdots②
\end{align*}

②から，
\begin{align*}
n(n-2)(S_{n+1}-S_n)&=S_n\\
n(n-2)S_{n+1}&=(n-1)^2S_n\quad\cdots③
\end{align*}

一方，①から$a_1=1$，②で$n=2$として$a_2=-1\quad\cdots④$である．$n\ge3$の時③から，
\begin{align*}
\frac n{n-1}S_{n+1}=\frac{n-1}{n-2}S_n=\cdots=\frac{3-1}{3-2}a_3
\end{align*}
\begin{align*}
\therefore\ S_n=\frac{n-2}{n-1}\cdot2a_3\quad(n\ge3)
\end{align*}

①から，$S_n=\dfrac{1-2/n}{1-1/n}\cdot2a_3\longrightarrow2a_3=1\quad\therefore\ a_3=\dfrac12$であるから，
\begin{align*}
S_n=\frac{n-2}{n-1}
\end{align*}

これは$n=2$でも成立（$\because④$）する．②から，$n\ge4$の時，
\begin{align*}
a_n=\frac1{(n-1)(n-3)}\cdot\frac{n-3}{n-2}=\frac1{(n-1)(n-2)}
\end{align*}

これは$n=3$で成立．以上から
\begin{align*}
a_n=\begin{cases}1&(n=1)\\-1&(n=2)\\\dfrac1{(n-1)(n-2)}&(n\ge3)\end{cases}
\end{align*}

\newpage

\end{document}
